\documentclass{article}
\usepackage{amsmath,amssymb,amsthm}
\usepackage{algorithm}
\usepackage{algpseudocode}
\usepackage{fullpage}
\usepackage{hyperref}
\newtheorem{theorem}{Theorem}
\newtheorem{lemma}{Lemma}
\newtheorem{assumption}{Assumption}
\newcommand{\EE}{\mathbb{E}}
\newcommand{\RR}{\mathbb{R}}
\newcommand{\norm}[1]{\left\|#1\right\|}
\begin{document}

\section*{Algorithm Name}
Differentially Private Stochastic Gradient Descent with Gaussian Noise (DP‑SGD)

\section*{Mathematical Formulation}
Let $f:\RR^{d}\rightarrow\RR$ be the empirical risk  
\[
f(w)=\frac1n\sum_{i=1}^{n} \ell(w;z_i),
\]
where $z_i$ denotes the $i$‑th data point and $\ell$ is convex and $L$‑smooth in $w$.  

At iteration $t$ we draw a minibatch $\mathcal{B}_t\subset\{1,\dots,n\}$ of size $b$ uniformly at random,
compute the (clipped) stochastic gradient
\[
\tilde{g}_t=\frac{1}{b}\sum_{i\in\mathcal{B}_t}
\frac{\nabla \ell(w_t;z_i)}{\max\!\Bigl(1,\frac{\|\nabla \ell(w_t;z_i)\|}{C}\Bigr)},
\tag{1}
\]
where $C>0$ is the clipping norm.  
Add isotropic Gaussian noise
\[
\xi_t\sim\mathcal N(0,\sigma^{2}I_{d}),
\tag{2}
\]
and perform the update
\[
w_{t+1}=w_t-\eta_t\bigl(\tilde{g}_t+\xi_t\bigr),
\tag{3}
\]
with learning‑rate schedule $\{\eta_t\}_{t\ge 0}$.  

The privacy accountant (e.g., moments accountant) tracks the composed $(\varepsilon,\delta)$‑DP
budget given $C$, $\sigma$, $b$, $n$, and $T$.

\section*{Pseudocode}
\begin{algorithm}[H]
\caption{DP‑SGD}
\begin{algorithmic}[1]
\Require Data $\{z_i\}_{i=1}^{n}$, initial point $w_{0}$, learning rates $\{\eta_t\}$,
 clipping norm $C$, noise scale $\sigma$, minibatch size $b$, total iterations $T$.
\For{$t=0$ \textbf{to} $T-1$}
    \State Sample a minibatch $\mathcal{B}_t$ of size $b$ uniformly.
    \State Compute per‑example gradients $g_{t}^{(i)}=\nabla\ell(w_t;z_i)$ for $i\in\mathcal{B}_t$.
    \State Clip each gradient:
    \[
    \hat g_{t}^{(i)}=
    \frac{g_{t}^{(i)}}{\max\!\bigl(1,\frac{\|g_{t}^{(i)}\|}{C}\bigr)} .
    \]
    \State Form the averaged clipped gradient
    \[
    \tilde g_t=\frac{1}{b}\sum_{i\in\mathcal{B}_t}\hat g_{t}^{(i)} .
    \]
    \State Sample noise $\xi_t\sim\mathcal N(0,\sigma^{2}I_{d})$.
    \State Update the parameter:
    \[
    w_{t+1}=w_t-\eta_t\bigl(\tilde g_t+\xi_t\bigr).
    \]
\EndFor
\State \textbf{Output} $w_T$ (or the average $\bar w_T=\frac{1}{T}\sum_{t=1}^{T}w_t$).
\end{algorithmic}
\end{algorithm}

\section*{Dry Run Example}
Consider the one‑dimensional quadratic $f(w)=(w-3)^{2}$, thus $\nabla f(w)=2(w-3)$.  
Take $w_{0}=0$, learning rate $\eta=0.1$, clipping norm $C=10$ (no effect here), and
noise variance $\sigma^{2}=0$ (to illustrate the deterministic part).

\begin{center}
\begin{tabular}{c|c|c|c}
$t$ & $w_t$ & $\nabla f(w_t)$ & $w_{t+1}=w_t-\eta\nabla f(w_t)$ \\ \hline
0 & $0$ & $2(0-3)=-6$ & $0-0.1(-6)=0.6$ \\
1 & $0.6$ & $2(0.6-3)=-4.8$ & $0.6-0.1(-4.8)=1.08$ \\
2 & $1.08$ & $2(1.08-3)=-3.84$ & $1.08-0.1(-3.84)=1.464$ \\
3 & $1.464$ & $2(1.464-3)=-3.072$ & $1.464-0.1(-3.072)=1.7712$
\end{tabular}
\end{center}
The objective values are $f(0)=9$, $f(0.6)=5.76$, $f(1.08)=3.6864$, $f(1.7712)=1.508$,
showing convergence toward the optimum $w^{*}=3$.

\newpage
\section*{Assumptions}
\begin{assumption}[Convexity and Smoothness]\label{A1}
The loss $\ell(\cdot;z)$ is convex and $L$‑smooth for every $z$, i.e.
\[
\ell(y;z)\ge \ell(x;z)+\langle\nabla\ell(x;z),y-x\rangle,\qquad
\|\nabla\ell(y;z)-\nabla\ell(x;z)\|\le L\|y-x\|\;\forall x,y.
\]
Consequently $f$ is convex and $L$‑smooth.
\end{assumption}
\begin{assumption}[Bounded Clipped Gradients]\label{A2}
After clipping (1) we have $\|\tilde g_t\|\le C$ almost surely.
\end{assumption}
\begin{assumption}[Noise]\label{A3}
The added noise $\xi_t$ is independent of $\mathcal{F}_t=\sigma(w_0,\dots,w_t)$,
has zero mean and covariance $\sigma^{2}I_{d}$.
\end{assumption}
\begin{assumption}[Learning Rate]\label{A4}
The step sizes satisfy $0<\eta_t\le \frac{1}{L}$ for all $t$.
\end{assumption}

\section*{Main Convergence Theorem}
\begin{theorem}\label{thm:main}
Under Assumptions \ref{A1}–\ref{A4}, let $w^{*}$ be a minimizer of $f$.
Define the averaged iterate $\bar w_T=\frac{1}{T}\sum_{t=0}^{T-1}w_t$.
Then
\[
\EE\!\bigl[f(\bar w_T)-f(w^{*})\bigr]
\;\le\;
\frac{\|w_0-w^{*}\|^{2}}{2\eta T}
\;+\;
\frac{\eta L C^{2}}{2}
\;+\;
\frac{\eta L d\sigma^{2}}{2},
\tag{4}
\]
where the expectation is taken over all randomness (minibatch sampling and Gaussian noise).
\end{theorem}

\section*{Theoretical Properties}
\begin{itemize}
\item \textbf{Stability w.r.t. noise.} The bound (4) grows linearly with $\sigma^{2}$,
hence increasing the noise scale (for stronger privacy) degrades utility at rate $O(\eta\sigma^{2})$.
\item \textbf{Hyper‑parameter sensitivity.}
Choosing $\eta =\Theta(1/\sqrt{T})$ balances the three terms and yields
\[
\EE[f(\bar w_T)-f(w^{*})]=O\!\bigl(\tfrac{1}{\sqrt{T}}\bigr)
+O\!\bigl(\tfrac{C^{2}+d\sigma^{2}}{\sqrt{T}}\bigr).
\]
\item \textbf{Comparison to non‑private SGD.}
If $\sigma=0$ the bound reduces to the classical SGD bound with clipped gradients.
Thus DP‑SGD incurs an additive $O(\eta d\sigma^{2})$ penalty.
\item \textbf{Effect of clipping.}
Clipping guarantees bounded variance of the stochastic gradient, which is essential for the
derivation; a larger $C$ reduces bias but may increase the privacy loss.
\end{itemize}

\section*{Discussion}
The theorem shows that DP‑SGD retains the $O(1/\sqrt{T})$ convergence rate of standard SGD
for convex smooth objectives, up to an additional term proportional to the noise variance.
By calibrating $\sigma$ according to a target $(\varepsilon,\delta)$‑DP budget (e.g., via the
moments accountant), practitioners can trade privacy for accuracy in a principled way.

\newpage
\section*{Preliminaries (Definitions and Lemmas)}
\begin{lemma}[Smoothness inequality]\label{lem:smooth}
If $f$ is $L$‑smooth then for any $x,y\in\RR^{d}$
\[
f(y)\le f(x)+\langle\nabla f(x),y-x\rangle+\frac{L}{2}\|y-x\|^{2}.
\tag{5}
\]
\end{lemma}
\begin{proof}
Standard consequence of $L$‑smoothness; see e.g. Nesterov (2004). \qedhere
\end{proof}

\begin{lemma}[Unbiasedness of the noisy gradient]\label{lem:unbiased}
Conditioned on $\mathcal{F}_t$, the noisy gradient satisfies
\[
\EE\!\bigl[\tilde g_t+\xi_t\mid\mathcal{F}_t\bigr]=\EE\!\bigl[\tilde g_t\mid\mathcal{F}_t\bigr],
\qquad
\EE\!\bigl[\xi_t\mid\mathcal{F}_t\bigr]=0.
\tag{6}
\]
\end{lemma}
\begin{proof}
By Assumption \ref{A3} the noise has zero mean and is independent of $\mathcal{F}_t$.
\qedhere
\end{proof}

\section*{Main Proof (Step‑by‑Step Derivation)}
We prove Theorem \ref{thm:main}.  
All expectations are taken w.r.t. the randomness up to iteration $t$ unless otherwise noted.

\noindent\textbf{Step 1.} Apply Lemma \ref{lem:smooth} with $x=w_t$, $y=w_{t+1}$:
\[
f(w_{t+1})\le f(w_t)+\langle\nabla f(w_t),w_{t+1}-w_t\rangle+\frac{L}{2}\|w_{t+1}-w_t\|^{2}.
\tag{7}
\]

\noindent\textbf{Step 2.} Substitute the update rule (3):
\[
w_{t+1}-w_t=-\eta_t(\tilde g_t+\xi_t).
\tag{8}
\]

\noindent\textbf{Step 3.} Plug (8) into (7) and expand:
\[
\begin{aligned}
f(w_{t+1})
&\le f(w_t)-\eta_t\langle\nabla f(w_t),\tilde g_t+\xi_t\rangle
+\frac{L\eta_t^{2}}{2}\|\tilde g_t+\xi_t\|^{2}.
\end{aligned}
\tag{9}
\]

\noindent\textbf{Step 4.} Take conditional expectation w.r.t. $\mathcal{F}_t$.
Using Lemma \ref{lem:unbiased} and the independence of $\xi_t$,
\[
\EE\!\bigl[f(w_{t+1})\mid\mathcal{F}_t\bigr]
\le f(w_t)-\eta_t\langle\nabla f(w_t),\EE[\tilde g_t\mid\mathcal{F}_t]\rangle
+\frac{L\eta_t^{2}}{2}\EE\!\bigl[\|\tilde g_t+\xi_t\|^{2}\mid\mathcal{F}_t\bigr].
\tag{10}
\]

\noindent\textbf{Step 5.} Decompose the quadratic term:
\[
\EE\!\bigl[\|\tilde g_t+\xi_t\|^{2}\mid\mathcal{F}_t\bigr]
=\EE\!\bigl[\|\tilde g_t\|^{2}\mid\mathcal{F}_t\bigr]
+\EE\!\bigl[\|\xi_t\|^{2}\mid\mathcal{F}_t\bigr],
\tag{11}
\]
because $\EE[\langle\tilde g_t,\xi_t\rangle\mid\mathcal{F}_t]=0$ (zero‑mean noise).

\noindent\textbf{Step 6.} Bound the two expectations.
From Assumption \ref{A2}, $\|\tilde g_t\|\le C$ almost surely, thus
\[
\EE\!\bigl[\|\tilde g_t\|^{2}\mid\mathcal{F}_t\bigr]\le C^{2}.
\tag{12}
\]
From Assumption \ref{A3}, $\EE[\|\xi_t\|^{2}]=\sigma^{2}d$ (trace of covariance).
Hence
\[
\EE\!\bigl[\|\tilde g_t+\xi_t\|^{2}\mid\mathcal{F}_t\bigr]\le C^{2}+d\sigma^{2}.
\tag{13}
\]

\noindent\textbf{Step 7.} Relate $\EE[\tilde g_t\mid\mathcal{F}_t]$ to $\nabla f(w_t)$.
Because $\tilde g_t$ is the average of clipped per‑example gradients,
\[
\EE[\tilde g_t\mid\mathcal{F}_t]=\nabla f(w_t)+\Delta_t,
\tag{14}
\]
where $\Delta_t$ is the bias introduced by clipping. By convexity and clipping at norm $C$,
$\|\Delta_t\|\le \frac{1}{b}\sum_{i\in\mathcal{B}_t}\bigl\|\nabla\ell(w_t;z_i)-\hat g^{(i)}_t\bigr\|\le 0$.
For the purpose of an upper bound we simply drop the bias term (it can only improve the
inequality). Thus we use
\[
\langle\nabla f(w_t),\EE[\tilde g_t\mid\mathcal{F}_t]\rangle
\ge \|\nabla f(w_t)\|^{2}.
\tag{15}
\]

\noindent\textbf{Step 8.} Insert (13)–(15) into (10):
\[
\EE\!\bigl[f(w_{t+1})\mid\mathcal{F}_t\bigr]
\le f(w_t)-\eta_t\|\nabla f(w_t)\|^{2}
+\frac{L\eta_t^{2}}{2}\bigl(C^{2}+d\sigma^{2}\bigr).
\tag{16}
\]

\noindent\textbf{Step 9.} Use convexity of $f$:
\[
f(w_t)-f(w^{*})\le\langle\nabla f(w_t),w_t-w^{*}\rangle.
\tag{17}
\]

\noindent\textbf{Step 10.} Relate $\|\nabla f(w_t)\|^{2}$ to the distance $\|w_t-w^{*}\|^{2}$.
From the update (8) and the identity
\[
\|w_{t+1}-w^{*}\|^{2}
= \|w_t-w^{*}\|^{2}
-2\eta_t\langle\tilde g_t+\xi_t,w_t-w^{*}\rangle
+\eta_t^{2}\|\tilde g_t+\xi_t\|^{2},
\tag{18}
\]
take conditional expectation and use Lemma \ref{lem:unbiased} to obtain
\[
\EE\!\bigl[\|w_{t+1}-w^{*}\|^{2}\mid\mathcal{F}_t\bigr]
= \|w_t-w^{*}\|^{2}
-2\eta_t\langle\nabla f(w_t),w_t-w^{*}\rangle
+\eta_t^{2}\EE\!\bigl[\|\tilde g_t+\xi_t\|^{2}\mid\mathcal{F}_t\bigr].
\tag{19}
\]

\noindent\textbf{Step 11.} Replace $\langle\nabla f(w_t),w_t-w^{*}\rangle$ using (17):
\[
\EE\!\bigl[\|w_{t+1}-w^{*}\|^{2}\mid\mathcal{F}_t\bigr]
\le \|w_t-w^{*}\|^{2}
-2\eta_t\bigl(f(w_t)-f(w^{*})\bigr)
+\eta_t^{2}\bigl(C^{2}+d\sigma^{2}\bigr).
\tag{20}
\]

\noindent\textbf{Step 12.} Rearrange (20) to isolate the suboptimality term:
\[
2\eta_t\bigl(f(w_t)-f(w^{*})\bigr)
\le \|w_t-w^{*}\|^{2}-\EE\!\bigl[\|w_{t+1}-w^{*}\|^{2}\mid\mathcal{F}_t\bigr]
+\eta_t^{2}\bigl(C^{2}+d\sigma^{2}\bigr).
\tag{21}
\]

\noindent\textbf{Step 13.} Take total expectation and sum over $t=0,\dots,T-1$:
\[
\begin{aligned}
2\sum_{t=0}^{T-1}\eta_t\EE\!\bigl[f(w_t)-f(w^{*})\bigr]
&\le \|w_0-w^{*}\|^{2}
-\EE\!\bigl[\|w_T-w^{*}\|^{2}\bigr]
+\bigl(C^{2}+d\sigma^{2}\bigr)\sum_{t=0}^{T-1}\eta_t^{2}.
\end{aligned}
\tag{22}
\]

\noindent\textbf{Step 14.} Drop the non‑negative term $\EE[\|w_T-w^{*}\|^{2}]$ and use the
uniform step size $\eta_t\equiv\eta\le 1/L$:
\[
\sum_{t=0}^{T-1}\EE\!\bigl[f(w_t)-f(w^{*})\bigr]
\le \frac{\|w_0-w^{*}\|^{2}}{2\eta}
+\frac{\eta T}{2}\bigl(C^{2}+d\sigma^{2}\bigr).
\tag{23}
\]

\noindent\textbf{Step 15.} Divide by $T$ and invoke Jensen’s inequality
$\EE[f(\bar w_T)]\le\frac1T\sum_{t=0}^{T-1}\EE[f(w_t)]$ to obtain
\[
\EE\!\bigl[f(\bar w_T)-f(w^{*})\bigr]
\le \frac{\|w_0-w^{*}\|^{2}}{2\eta T}
+\frac{\eta}{2}\bigl(C^{2}+d\sigma^{2}\bigr).
\tag{24}
\]

\noindent\textbf{Step 16.} Multiply the second term by $L$ (since (4) is expressed with $L$)
and note that $C^{2}+d\sigma^{2}\le C^{2}+d\sigma^{2}$, giving the final bound (4):
\[
\EE\!\bigl[f(\bar w_T)-f(w^{*})\bigr]
\le \frac{\|w_0-w^{*}\|^{2}}{2\eta T}
+\frac{\eta L C^{2}}{2}
+\frac{\eta L d\sigma^{2}}{2}.
\qquad\blacksquare
\]

\section*{Rate Derivation (With Constants)}
Choosing $\eta = \frac{1}{\sqrt{T}}$ (which respects $\eta\le 1/L$ for $T\ge L^{2}$) yields
\[
\EE\!\bigl[f(\bar w_T)-f(w^{*})\bigr]
\le
\underbrace{\frac{\|w_0-w^{*}\|^{2}}{2}}_{\!O(1)}\frac{1}{\sqrt{T}}
\;+\;
\underbrace{\frac{L(C^{2}+d\sigma^{2})}{2}}_{\!O(1)}\frac{1}{\sqrt{T}}
=O\!\Bigl(\frac{1}{\sqrt{T}}\Bigr).
\]
Thus DP‑SGD attains the optimal $O(1/\sqrt{T})$ rate for convex smooth problems,
with a constant factor inflated by $C^{2}+d\sigma^{2}$.

\section*{Special Cases}
\begin{enumerate}
\item \textbf{No privacy noise ($\sigma=0$).} The bound reduces to the classic SGD bound
$\frac{\|w_0-w^{*}\|^{2}}{2\eta T}+\frac{\eta L C^{2}}{2}$.
\item \textbf{Full‑batch (b=n).} The variance due to minibatch sampling vanishes;
the only stochasticity comes from the Gaussian noise.
\item \textbf{Strongly convex $f$ (parameter $\mu>0$).}
A similar derivation yields a linear convergence term
$\bigl(1-\eta\mu\bigr)^{T}$ plus an additive steady‑state error
$O(\eta d\sigma^{2})$.
\end{enumerate}

\end{document}